\section{Conclusions and Future Work}
\label{sec:conclusions}

This work studied provincial differences in female employment in Italy (year 2022) and their relationship with ECEC availability and socio-economic covariates, using a Bayesian spatial modeling framework with a CAR random effect and a sparse finite mixture approach for model-based clustering.\\

The label-assignment step resulted in a single cluster including all provinces, indicating no posterior support for distinct groups characterized by different regression effects. Consequently, inference is effectively that of a global spatial regression: associations between the selected covariates and female employment are summarized through one coefficient vector, while remaining geographic disparities are
captured by the spatial random effect.\\

Future work could extend the analysis in several directions, one option could be repeating the study over multiple years to increase statistical power to detect heterogeneous relationships and enable joint spatio-temporal modeling. Another option would be to enrich the dataset with additional covariates (not correlated with the others we already have) that capture further drivers of the outcome, thereby improving explanatory power and reducing the share of variation absorbed by the spatial effect.\\


% REMEMBER TO REMOVE NOCITE{*} FROM THE BIB, OTHERWISE THERE ARE PAPERS WE DIDN'T USE IN THE BIB